\documentclass[11pt]{article}
\usepackage[utf8]{inputenc}
\usepackage[T1]{fontenc}
\usepackage{lmodern}
\usepackage{amsmath,amssymb,amsthm}
\usepackage[margin=1in]{geometry}
\usepackage{booktabs}
\usepackage{tabularx}
\usepackage{listings}
\usepackage{xcolor}
\usepackage[hidelinks]{hyperref}
\usepackage{url}

\newtheorem{theorem}{Theorem}
\newtheorem{corollary}[theorem]{Corollary}
\newtheorem{lemma}[theorem]{Lemma}
\newtheorem{proposition}[theorem]{Proposition}
\theoremstyle{definition}
\newtheorem{definition}[theorem]{Definition}
\theoremstyle{remark}
\newtheorem{remark}[theorem]{Remark}

\newcommand{\lean}[1]{\texttt{\detokenize{#1}}}
\setlength{\emergencystretch}{3em}
\newcommand{\Rea}{\operatorname{Re}}
\newcommand{\Ima}{\operatorname{Im}}

\lstdefinelanguage{lean}{
  morekeywords={theorem,def,noncomputable,lemma,by,fun,namespace,end,open,import,Prop,Set},
  sensitive=true, morecomment=[l]{--}}
\lstset{language=lean, basicstyle=\ttfamily\small, columns=fullflexible, keepspaces=true,
  extendedchars=true, breaklines=true, frame=single, framesep=4pt, xleftmargin=6pt, xrightmargin=6pt,
  keywordstyle=\bfseries,
  literate=
    {∀}{{$\forall$}}1 {∃}{{$\exists$}}1 {∈}{{$\in$}}1 {≤}{{$\le$}}1 {≥}{{$\ge$}}1
    {→}{{$\to$}}1 {↔}{{$\leftrightarrow$}}1 {ρ}{{$\rho$}}1 {σ}{{$\sigma$}}1 {φ}{{$\varphi$}}1
    {θ}{{$\theta$}}1 {ℂ}{{$\mathbb{C}$}}1 {ℝ}{{$\mathbb{R}$}}1 {ℕ}{{$\mathbb{N}$}}1
    {≠}{{$\ne$}}1 {∧}{{$\wedge$}}1 {∑}{{$\Sigma$}}1 {₀}{{$_0$}}1 {₁}{{$_1$}}1 {₂}{{$_2$}}1
    {π}{{$\pi$}}1}

\title{The Riemann hypothesis up to height $634.5$,\\ checked by the Lean kernel}
\author{Samuel Lavery}
\date{Draft, 29 September 2026}

\begin{document}
\maketitle

\begin{abstract}
We give a proof, checked by the kernel of the Lean~4 proof assistant against Mathlib's definition of the Riemann zeta
function, that every nontrivial zero $\rho$ of $\zeta$ with $|\Ima\rho|\le 634.5$ lies on the critical line
$\Rea s=\tfrac12$, and that every such zero with $0<\Ima\rho\le 634.5$ is a simple zero of the Riemann $\xi$-function.
A multiplicity-weighted count of $366$ zeros in the window is also checked, so there are exactly $366$ zeros with
$0<\Ima\rho\le634.5$. The proof follows Turing's scheme. Sign changes of the Hardy $Z$-function certify zeros on the
line, and an argument-principle count certifies that there are no others. Every numerical fact enters as a rational
inequality whose proof the kernel checks. The only axioms used are Lean's three standard ones, and the development
does not use native code evaluation. We also derive the instance $T=634.5$ of the predicate \lean{riemannZeta.RH_up_to}
of the PrimeNumberTheorem+ project, which at the revision we use records the numerical verifications of Platt, of
Gourdon--Wedeniwski and of Platt--Trudgian as statements without proofs. The height is small by the standards of
numerical verification. The contribution is a verification whose trusted base is the Lean kernel and the definitions,
not numerical software. The development consists of $367$ generated certificate modules ($546$\,MB of Lean source), and it
builds on a laptop under a documented memory limit.
\end{abstract}

\section{Introduction}\label{sec:intro}

Numerical verification of the Riemann hypothesis (RH) up to a height $T$ means proving that every zero $\beta+i\gamma$
of $\zeta$ with $0<\gamma\le T$ has $\beta=\tfrac12$. The method in use since Turing~\cite{Turing1953} has two parts.
Zeros on the critical line are located by sign changes of the real-valued Hardy function
$Z(t)=e^{i\theta(t)}\zeta(\tfrac12+it)$. An independent count of all zeros up to height $T$ then shows that the list is
complete. The current record is due to Platt and Trudgian~\cite{PlattTrudgian2021}, who ``verify numerically, in a
rigorous way using interval arithmetic, that the Riemann hypothesis is true up to height $3\cdot10^{12}$''.

These computations are rigorous in the sense of interval arithmetic, and they run in ordinary software. Formal
libraries use their conclusions as hypotheses. At revision \lean{21998bb} the PrimeNumberTheorem+ (PNT+)
project~\cite{PNTplus} states three such results in Lean, each closed by \lean{sorry}:
\begin{lstlisting}
theorem Platt_theorem : riemannZeta.RH_up_to 30610046000 := sorry
theorem GW_theorem    : riemannZeta.RH_up_to 2445999556030 := sorry
theorem PT_theorem_1  : riemannZeta.RH_up_to 3e12 := sorry
\end{lstlisting}
This paper gives a kernel-checked proof of the same kind of statement at a much smaller height.

\paragraph{Results.} All theorems below are checked by the Lean kernel, version 4.34.0, with Mathlib at tag
\lean{v4.34.0}. Their axiom footprint is exactly \lean{[propext, Classical.choice, Quot.sound]}.
\begin{enumerate}
\item Every nontrivial zero with $|\Ima\rho|\le 634.5$ has $\Rea\rho=\tfrac12$ (Theorem~\ref{thm:window}).
\item Every nontrivial zero with $0<\Ima\rho\le 634.5$ is a simple zero of $\xi$ (Theorem~\ref{thm:simple}), and the
  multiplicity-weighted count in that window is $366$ (Theorem~\ref{thm:census}).
\item PNT+'s \lean{riemannZeta.RH_up_to (1269/2)} holds (Theorem~\ref{thm:pnt}).
\item Mathlib's \lean{RiemannHypothesis} follows from one explicitly named hypothesis about the paired Dirichlet eta
  series: at each height it vanishes for at most one abscissa in the strip (Section~\ref{sec:interface}). The
  hypothesis itself is not proved here beyond height $634.5$.
\end{enumerate}

\paragraph{Scope.} Nothing here is new as mathematics. Zeros of $\zeta$ in this range were computed, and shown to lie
on the critical line, in the 1930s; see~\cite{Titchmarsh} for the history. $634.5$ is also tiny next to $3\cdot10^{12}$:
the window holds $366$ zeros, against about $1.2\cdot10^{13}$. It lies far below every height that PNT+ records as a
theorem without proof ($3.06\cdot10^{10}$ and above). The one contribution is the trusted base of a classical
computation. That base shrinks to the Lean kernel and to Mathlib's
definitions of $\zeta$, $\Gamma$, $\log$, $\cos$ and $\pi$. The numerical data enters only as rational witnesses that
the kernel re-checks. Section~\ref{sec:scaling} discusses how the cost grows with height.

\section{Statements and trust base}\label{sec:statements}

\subsection{Definitions}
Mathlib defines \lean{riemannZeta} on all of $\mathbb C$. At the pole $s=1$ it takes a fixed value that plays no role
here. The nontrivial zeros are
\begin{lstlisting}
def ZD.NontrivialZeros : Set ℂ :=
  { s : ℂ | 0 < s.re ∧ s.re < 1 ∧ riemannZeta s = 0 }
\end{lstlisting}
Mathlib's \lean{completedRiemannZeta} is $\Lambda(s)=\pi^{-s/2}\Gamma(s/2)\zeta(s)$. The development uses the Riemann
$\xi$-function in the form
\[
  \xi(s) \;=\; \tfrac{s(s-1)}{2}\,\Lambda_0(s)+\tfrac12 \;=\; \tfrac12 s(s-1)\,\pi^{-s/2}\Gamma(s/2)\zeta(s),
\]
where $\Lambda_0$ is Mathlib's entire part \texttt{completedRiemannZeta}$_0$. The multiplicity of a zero is the order of
vanishing of $\xi$, \lean{ZD.xiOrderNat z := analyticOrderNatAt riemannXi z}.

\subsection{Main theorems}

\begin{theorem}[the window]\label{thm:window}
Every nontrivial zero $\rho$ of $\zeta$ with $|\Ima\rho|\le 634.5$ satisfies $\Rea\rho=\tfrac12$.
\begin{lstlisting}
theorem CriticalLinePhasor.DVP.online_below_1269o2 :
    ∀ ρ ∈ ZD.NontrivialZeros, |ρ.im| ≤ 1269/2 → ρ.re = 1/2
\end{lstlisting}
\end{theorem}

\begin{theorem}[the census]\label{thm:census}
The number of zeros with $0<\Rea s<1$ and $0<\Ima s\le 634.5$, counted with multiplicity, is $366$.
\begin{lstlisting}
def stripZeroCountMult (T : ℝ) : ℕ :=
  ∑ ρ ∈ stripZeroFinset T, ZD.xiOrderNat ρ
theorem CriticalLinePhasor.DVP.stripZeroCountMult_1269o2 :
    stripZeroCountMult (1269/2) = 366
\end{lstlisting}
Here \lean{stripZeroFinset T} is the finite set of nontrivial zeros with $0<\Ima\rho\le T$.
\end{theorem}

\begin{theorem}[simplicity]\label{thm:simple}
Every nontrivial zero $\rho$ with $0<\Ima\rho\le 634.5$ is a simple zero of $\xi$.
\begin{lstlisting}
theorem VerifiedWindow.xiOrder_eq_one_below_1269o2 {ρ : ℂ}
    (hρ : ρ ∈ ZD.NontrivialZeros) (hpos : 0 < ρ.im) (hle : ρ.im ≤ 1269 / 2) :
    ZD.xiOrderNat ρ = 1
\end{lstlisting}
\end{theorem}
Theorems~\ref{thm:census} and~\ref{thm:simple} together say there are exactly $366$ distinct zeros with
$0<\Ima\rho\le634.5$. In the open strip, $\pi^{-s/2}\Gamma(s/2)$ and $\tfrac12 s(s-1)$ have no zeros, so these are also
simple zeros of $\zeta$. We state this in prose; the compiled statement is about $\xi$.

\begin{theorem}[PNT+'s predicate]\label{thm:pnt}
With PNT+'s definitions (file \lean{IEANTN/ZetaDefinitions.lean} at revision \lean{21998bb})
\begin{lstlisting}
noncomputable def riemannZeta.zeroes : Set ℂ := {s : ℂ | riemannZeta s = 0}
noncomputable def riemannZeta.zeroes_rect (I J : Set ℝ) : Set ℂ :=
  {s : ℂ | s.re ∈ I ∧ s.im ∈ J ∧ s ∈ zeroes}
noncomputable def riemannZeta.RH_up_to (T : ℝ) : Prop :=
  IsEmpty (riemannZeta.zeroes_rect (Set.Ioo 0.5 1) (Set.Icc 0 T))
\end{lstlisting}
we have
\begin{lstlisting}
theorem VerifiedWindow.pnt_RH_up_to_1269o2 : riemannZeta.RH_up_to (1269 / 2)
\end{lstlisting}
\end{theorem}
The proof derives it from Theorem~\ref{thm:window}. A zero in PNT+'s rectangle is a nontrivial zero with
$|\Ima s|\le 634.5$ and $\Rea s>\tfrac12$.

\subsection{What is trusted}
\begin{itemize}
\item \textbf{Trusted:} the Lean~4 kernel (v4.34.0). The Mathlib definitions of $\zeta$, $\Gamma$, $\log$, $\cos$,
  $\sin$, $\pi$ and analytic order at tag \lean{v4.34.0}. For Theorem~\ref{thm:pnt} only, the three PNT+ definitions
  displayed above.
\item \textbf{Axioms:} \texttt{\#print axioms} reports exactly \lean{[propext, Classical.choice, Quot.sound]} for every
  theorem in this section. The dependency closure contains no \lean{sorry} and no user axiom. None of the ladder,
  station, census or segment modules uses \lean{native_decide}, so \lean{Lean.ofReduceBool} does not appear.
\item \textbf{Untrusted:} the Python programs that generate the certificate modules. They choose the station heights,
  the rational phase stand-ins and the rational bounds. A wrong choice makes a module fail to compile. It cannot make a
  false statement compile.
\end{itemize}

\section{Method}\label{sec:method}

The proof is Turing's scheme carried out inside the kernel. We describe it in standard terms; Table~\ref{tab:names} maps
the Lean names.

\subsection{Hardy's function and the sign at a station}
On the critical line, $\Lambda(\tfrac12+it)$ is real. The development works with
\[
  G(t) \;=\; \Rea\Lambda(\tfrac12+it) \;=\; e^{\Rea\log\Gamma_{\mathbb R}(\frac12+it)}\,Z(t),
  \qquad Z(t)=e^{i\theta(t)}\zeta(\tfrac12+it),
\]
where $\theta(t)=\int_0^t\Rea\,(\Gamma_{\mathbb R}'/\Gamma_{\mathbb R})(\tfrac12+iu)\,du$ is the Riemann--Siegel
theta function. $G$ and $Z$ have the same sign. A sign change of $G$ on $[a,b]$ gives a zero on the line in $(a,b)$ by the
intermediate value theorem (\lean{hardyG_sign_change1269o2}).

The sign is read through a \emph{rational proxy phase}. Suppose $|\varphi-\theta(t)|<\pi/2$. Then a positive value of
$\cos\varphi\,\Rea\zeta(\tfrac12+it)-\sin\varphi\,\Ima\zeta(\tfrac12+it)$ implies $Z(t)>0$, and a negative value implies
$Z(t)<0$ (\lean{hardyG_sign_of_proxy}). The tolerance is more than a radian. So $\theta(t)$ is never needed to many digits. The
development proves an enclosure $|\theta(30)+\tfrac{15}{2}\pi-31.6323|\le 0.105$ (\lean{theta30_enclosure}) and then
the increment bound
\[
  \bigl|(\theta(T)-\theta(T_1))-(\theta_{\mathrm{main}}(T)-\theta_{\mathrm{main}}(T_1))\bigr|\le 10\Bigl(\frac1{T_1}-\frac1T\Bigr),
  \qquad \theta_{\mathrm{main}}(T)=\tfrac T2\log\tfrac{T}{2\pi}-\tfrac T2
\]
(\lean{theta_increment_enclosure_from}). Together these place a rational $\varphi$ plus an integer multiple of $2\pi$
within $\pi/2$ of $\theta(t)$ at every station.

\subsection{Certified evaluation}
At a station $t$ the quantity $\Rea\bigl(e^{-i\varphi}\zeta(\tfrac12+it)\bigr)$ is bounded below or above using the
Euler--Maclaurin formula truncated after $N$ terms. $N$ is chosen per station, between $200$ and $4000$; most stations use
$N=250$ or $N=900$. The boundary term $N^{1-s}/(s-1)$ is evaluated exactly (\lean{boundary_re_eval_half}), and the
remainder is bounded explicitly. Each summand $n^{-1/2}\cos(t\log n-\varphi)$ is enclosed by rationals from three ingredients:
\begin{itemize}
\item tables of rational brackets $\ell_n^-\le\log n\le\ell_n^+$ for $n\le 4000$, each proved from bounds on the
  exponential;
\item rational brackets for $n^{-1/2}$;
\item Taylor polynomials for $\cos$ and $\sin$ with explicit remainder.
\end{itemize}
The station's conclusion is a single inequality between rational numbers, checked by \lean{norm_num}. A station is
therefore a large but elementary certificate. At $t=150.5$ it runs to $22{,}162$ lines.

\subsection{The count}
Completeness comes from the argument principle. Call a height $T$ \emph{good} when a certificate shows
$\Rea\zeta(\sigma+iT)\ge\tfrac15$ for $\tfrac12\le\sigma\le2$ (\lean{re_pos_at_1269o2}). That certificate comes from an
Euler--Maclaurin lower bound with $N=1000$ on seven $\sigma$-intervals (\lean{re_zeta1269o2_lower}). At a good height,
the continuous logarithm of $\zeta$ along the top of the standard contour stays on the principal branch. This gives
$|S(T)|=|\arg\zeta(\tfrac12+iT)|/\pi<\tfrac12$. The compiled argument-principle identity
(\lean{classicalSContour_eq_Smult_add_offLine}) reads $N(T)=\theta(T)/\pi+1+S(T)$ with $N(T)$ counted with multiplicity.
Combined with the enclosure of $\theta(634.5)$, it pins the integer $N(634.5)$ to $366$ (Theorem~\ref{thm:census}).

\subsection{Exhaustion and simplicity}
The ladder is $367$ stations $10\le t_0<t_1<\dots<t_{366}\le634.5$ with $G(t_i)G(t_{i+1})<0$ for every $i$
(\lean{LadderData1269o2}, discharged by \lean{ladder1269o2Data_holds}). It gives $366$ distinct zeros on the line. Any
further zero in the window, on or off the line, or any zero of order at least two, would push the weighted count above
$366$. The same finite-set argument proves Theorem~\ref{thm:window} (\lean{online_below_1269o2_of_ladder}) and
Theorem~\ref{thm:simple}. Zeros with negative imaginary part follow by conjugation.

\subsection{Rungs}
The window was built in $31$ rungs, with top heights $58, 63, 74, 81, 91, 109, 126, 144.5, \dots, 607.5, 634.5$. Each
rung has its own good height, segment certificate, census and exhaustion. The final theorems use only the top rung. The
bottom rung, \lean{online_below_58}, is the smallest slice. It uses $13$ stations and is suitable for a quick check
(Section~\ref{sec:engineering}).

\begin{table}[t]
\centering\small
\begin{tabularx}{\linewidth}{@{}lX@{}}
\toprule
Lean name & Standard meaning \\
\midrule
\lean{riemannZeta}, \lean{completedRiemannZeta} & $\zeta(s)$; $\Lambda(s)=\pi^{-s/2}\Gamma(s/2)\zeta(s)$ \\
\lean{ZD.riemannXi}, \lean{ZD.xiOrderNat} & $\xi(s)=\tfrac12s(s-1)\Lambda(s)$; order of vanishing of $\xi$ \\
\lean{ZD.NontrivialZeros} & zeros with $0<\Rea s<1$ \\
\lean{line t}, \lean{theta} & $\tfrac12+it$; Riemann--Siegel $\theta(t)$ \\
\lean{hardyG} & $\Rea\Lambda(\tfrac12+it)$, a positive multiple of Hardy's $Z(t)$ \\
\lean{hardyG_sign_of_proxy} & sign of $Z$ through a phase within $\pi/2$ of $\theta(t)$ \\
\lean{station_301o2_sign} (etc.) & certified sign of $Z$ at one station \\
\lean{LadderData1269o2} & $367$ stations with alternating signs \\
\lean{GoodHeight} & no zeros on the top edge of the contour at height $T$ \\
\lean{classicalSContour} & $S(T)$ computed along the contour \\
\lean{stripZeroCountMult} & $N(T)$, counted with multiplicity \\
\lean{laneChart3D} & the paired eta series $\sum_k((2k+1)^{-s}-(2k+2)^{-s})$ at $s=\sigma+iy$ \\
\lean{OneVanishingPerHeight3D} & for each $y$, at most one $\sigma\in(0,1)$ with $\eta(\sigma+iy)=0$ \\
\bottomrule
\end{tabularx}
\caption{Lean names and their standard meanings. The Lean identifiers come from a larger development and keep its
naming.}\label{tab:names}
\end{table}

\section{The one-hypothesis interface}\label{sec:interface}

Let $\eta(s)=\sum_{n\ge1}(-1)^{n-1}n^{-s}=(1-2^{1-s})\zeta(s)$ for $\Rea s>0$. The development represents the paired
series $\sum_k\bigl((2k+1)^{-s}-(2k+2)^{-s}\bigr)$ as a sum of planar vectors, \lean{laneChart3D}~$\sigma$~$y$ at
$s=\sigma+iy$. It proves that this vanishes for $0<\sigma<1$ exactly when $\zeta(\sigma+iy)=0$
(\lean{fiber3D_vanishes_iff_zero}). The factor $1-2^{1-s}$ vanishes only on $\Rea s=1$. The named hypothesis is
\begin{lstlisting}
def CarrierReadout.OneVanishingPerHeight3D : Prop :=
  ∀ (y σ₁ σ₂ : ℝ), 0 < σ₁ → σ₁ < 1 → 0 < σ₂ → σ₂ < 1 →
    laneChart3D σ₁ y = 0 → laneChart3D σ₂ y = 0 → σ₁ = σ₂
\end{lstlisting}
and the compiled theorem is
\begin{lstlisting}
theorem OneAssumptionRH.riemannHypothesis_of_oneAssumption
    (h : CarrierReadout.OneVanishingPerHeight3D) : RiemannHypothesis
\end{lstlisting}
where \lean{RiemannHypothesis} is Mathlib's statement. The proof is short. By the functional equation and conjugation,
a zero at $\sigma+iy$ comes with a zero at $(1-\sigma)+iy$ (\lean{fiber3D_vanishes_at_partner}). The hypothesis then
forces $\sigma=1-\sigma$. The bridge \lean{RHBridge.no_offline_zeros_implies_rh} converts the result to Mathlib's form.
The hypothesis restricted to $|y|\le634.5$ is a theorem (\lean{CarrierReadout.oneVanishingPerHeight3D_below}), by
Theorem~\ref{thm:window}. Above that height it is not proved here.

\section{Engineering and reproducibility}\label{sec:engineering}

\paragraph{Size.} The window rests on $367$ generated station modules totalling $546.2$\,MB of Lean source. Around them
sit $31$ rungs of census, segment and exhaustion modules, the rational log and square-root tables to $4000$, and the
contour-integral machinery. A Python program (\lean{tmp/att271_station_emitter.py} in the repository) emits the station
modules from witness tables. It is untrusted.

\paragraph{Throttling.} The station modules dominate the cost. At default parallelism, which here is one module per
core on $12$ cores, the build exhausted a $32$\,GB laptop. Every build reported here used
\begin{lstlisting}
env LEAN_NUM_THREADS=2 nice -n 10 lake build <Module>
\end{lstlisting}
which caps Lake at about two concurrent modules.

\paragraph{Measured costs} (Apple M2 Max, 12 cores, 32\,GB, macOS; wall time and peak resident memory from
\lean{/usr/bin/time -l}):
\begin{center}\small
\begin{tabular}{@{}lrr@{}}
\toprule
Task & Wall time & Peak RSS \\
\midrule
PNT+ \lean{IEANTN.ZetaDefinitions} (with its PNT+ dependency) & $35.6$\,s & $3.7$\,GB \\
Load full window and check the bridge file & $282.6$\,s & $5.9$\,GB \\
Load full window, print axioms only & $280.3$\,s & $6.2$\,GB \\
Re-check a station module from source ($60$ modules, heights $150$--$256$) & $51$--$247$\,s (median $75$) & --- \\
Fast slice: re-check the 15 height-58 modules from source$^{\dagger}$ & $741.6$\,s & $6.5$\,GB \\
Write the bridge module's olean (\lean{lake env lean -o}) & $257.8$\,s & $9.1$\,GB \\
Run \lean{Check.lean} (all statements and axiom footprints) & $251.3$\,s & $8.8$\,GB \\
\bottomrule
\end{tabular}

\smallskip
\footnotesize $^{\dagger}$One \lean{lake env lean} at a time against existing oleans, while a separate two-module build was
running; the times are upper bounds for an idle machine. Peak memory is dominated by loading Mathlib.
\end{center}

\paragraph{Checking.} In the directory \lean{papers/rh_window/}, the file \lean{Check.lean} imports the capstones. It prints
the statement and axiom footprint of every theorem in Section~\ref{sec:statements} and Section~\ref{sec:interface}. The
file \lean{ARTIFACT.md} in the same directory gives the full build recipe.

\paragraph{Build state.} In the working tree used for this draft, Lake's traces marked some station oleans as out of date
with respect to their sources. Before we stopped that rebuild, $60$ of them had been re-checked from the current sources
without error. Importing an olean does not re-run the kernel on its proofs. A reader should therefore certify the
artifact by one of two routes: a complete build from a clean checkout, or a replay of the kernel over the oleans with
\lean{lean4checker}. Neither has been run for this draft.

\section{Scaling with height}\label{sec:scaling}

Two things grow with $T$: the number of stations and the size of each. The number of zeros up to $T$ is
$N(T)\approx\frac{T}{2\pi}\log\frac{T}{2\pi e}+\frac78$, which is about $366$ at $634.5$, $415$ at $700$, $649$ at
$1000$ and $1517$ at $2000$. The ladder uses one station per zero, plus one. Station modules grow with height as the
Euler--Maclaurin truncation lengthens. A least-squares fit over the $367$ modules with $t\ge50$ gives size roughly
proportional to $t^{1.26}$, with medians of $0.17$\,MB below $t=100$ and $1.84$\,MB between $500$ and $634.5$. Extending to
$T=1000$ would add about $283$ stations of $2$--$4$\,MB each. At the per-station re-check times measured above, that is
days of throttled laptop time, not months. These figures are extrapolations, not measurements. Very large heights
would need a different evaluation, such as the Riemann--Siegel formula or the Odlyzko--Sch\"onhage algorithm, in place
of Euler--Maclaurin, and certified arithmetic well beyond hand-written rational tables.

\section{Related work}\label{sec:related}

\paragraph{Numerical verification.} Platt and Trudgian~\cite{PlattTrudgian2021}: ``We verify numerically, in a rigorous way
using interval arithmetic, that the Riemann hypothesis is true up to height $3\cdot10^{12}$.'' The methods descend from
Turing~\cite{Turing1953}, with corrections by Lehman~\cite{Lehman1970}; see Booker~\cite{Booker2006} for an account.

\paragraph{Formal libraries.} PNT+~\cite{PNTplus} defines \lean{riemannZeta.RH_up_to} and states the verifications of
Platt, of Gourdon--Wedeniwski and of Platt--Trudgian with \lean{sorry} (Section~\ref{sec:intro}). On the Lean community
Zulip, in January 2026, the PNT+ discussion ``N(T)'' asked whether Mathlib or PNT+ had machinery to define the zero
count $N(T)$. The reply was that it did not yet exist there. The definitions \lean{zeroes} and \lean{zeroes_rect} used in
Theorem~\ref{thm:pnt} were proposed in that discussion. Loeffler and Stoll~\cite{LoefflerStoll}
formalize $\zeta$ and Dirichlet $L$-functions in Mathlib, including a formal statement of RH. In Isabelle, Eberl's
entry~\cite{EberlAFP} ``builds upon the results about formal and analytic Dirichlet series to define the Hurwitz
$\zeta$ function $\zeta(a,s)$ and, based on that, the Riemann $\zeta$ function $\zeta(s)$''. We found no entry in the
Archive of Formal Proofs that verifies zeros of $\zeta$.

\paragraph{Zeros in proof assistants.} Alpöge and Furman~\cite{AlpogeFurman} prove that ``at least two thirds of the
nontrivial zeros of the Riemann zeta function, counted with multiplicity, are simple and lie on the critical line''.
They provide a Lean~4 formalization that includes the Riemann--von Mangoldt counting formulas. That result concerns
proportions of zeros at large height, not the location of individual zeros. A recent upload by Islam~\cite{Islam2026}
states that ``the Real part is proved, for $\zeta$ to height $3\times10^{12}$ on the certificate of Platt and
Trudgian''. It takes the numerical certificate as an input rather than checking it in the kernel. Several public
repositories formalize Hardy's theorem or zero counting in Lean without placing individual zeros.

\paragraph{Priority.} We searched the sources above, the PNT+ discussion on the Lean Zulip and the Archive of Formal
Proofs. We found no earlier proof-assistant verification of RH up to any positive height. The search is not exhaustive.
We make no stronger claim.

\section*{Disclosure of AI assistance}
\emph{[For the author to finalize.]} The Lean development and this manuscript were produced with extensive assistance
from Claude, an AI system developed by Anthropic. It helped generate the certificate modules, write the proofs of the
bridge and simplicity theorems, and draft the text. Every theorem stated here has been checked by the Lean kernel as
described in Section~\ref{sec:statements}.

\begin{thebibliography}{99}
\bibitem{AlpogeFurman} L.~Alpöge and R.~Furman, \emph{More than two thirds of the zeta zeros are simple and on the
  critical line}, arXiv:2608.13637, 2026. Lean formalization: \url{https://github.com/anthropics/zeta-23-lean}.
\bibitem{Booker2006} A.~R.~Booker, \emph{Turing and the Riemann hypothesis}, Notices Amer. Math. Soc. \textbf{53}
  (2006), no.~10, 1208--1211.
\bibitem{EberlAFP} M.~Eberl, \emph{The Hurwitz and Riemann $\zeta$ functions}, Archive of Formal Proofs, October 2017.
\bibitem{Islam2026} M.~F.~Islam, \emph{A formal proof of the real part of the Riemann hypothesis}, Zenodo record
  22949645, September 2026.
\bibitem{Lehman1970} R.~S.~Lehman, \emph{On the distribution of zeros of the Riemann zeta-function}, Proc. London Math.
  Soc. (3) \textbf{20} (1970), 303--320.
\bibitem{LoefflerStoll} D.~Loeffler and M.~Stoll, \emph{Formalizing zeta and $L$-functions in Lean}, Annals of Formalized
  Mathematics; arXiv:2503.00959, 2025.
\bibitem{PlattTrudgian2021} D.~Platt and T.~Trudgian, \emph{The Riemann hypothesis is true up to $3\cdot10^{12}$}, Bull.
  London Math. Soc. \textbf{53} (2021), no.~3, 792--797.
\bibitem{PNTplus} A.~Kontorovich, T.~Tao et al., \emph{PrimeNumberTheorem+},
  \url{https://github.com/AlexKontorovich/PrimeNumberTheoremAnd}, revision \texttt{21998bb}.
\bibitem{Titchmarsh} E.~C.~Titchmarsh, \emph{The Theory of the Riemann Zeta-Function}, 2nd ed., revised by
  D.~R.~Heath-Brown, Oxford University Press, 1986.
\bibitem{Turing1953} A.~M.~Turing, \emph{Some calculations of the Riemann zeta-function}, Proc. London Math. Soc. (3)
  \textbf{3} (1953), 99--117.
\end{thebibliography}

\end{document}
